\section{The fixed-shift affine M\"obius interface}
\label{sec:affine-interface}

The preceding sections show that neither Fourier centering,
conditional expectation, nor content inversion removes \(Z_X\).
Opening an ultra divisor on one polarized leg instead exposes a
two-sign affine M\"obius correlation.  We record the exact form and
the costs that remain outside it.

\subsection{Reflection of one physical column}

Fix one opposite row and one of the two polarizations in
\eqref{eq:literal-two-leg-polarization}.  On the opened ultra leg,
write a row as \(m=\ell d\), and reflect a divisor \(u\) by
\begin{equation}
 \ell d j+h_0=us.
 \label{eq:affine-reflection}
\end{equation}
The actual row-gcd mask can be expanded after the opposite row has
been fixed.  On one divisor layer \(v\), write
\[
 d=v\widetilde d.
\]
Primitive support gives
\begin{equation}
 (s,\ell jv)=1,
 \qquad
 (s\ell jv,h_0)=1.
 \label{eq:affine-slope-coprimality}
\end{equation}
Hence there is a unique residue
\(\widetilde d_0\pmod s\) satisfying
\[
 \ell jv\widetilde d_0+h_0\equiv0\pmod s.
\]
Put
\begin{equation}
 \widetilde u_0
 =\frac{\ell jv\widetilde d_0+h_0}{s},
 \qquad
 a=\ell jv.
 \label{eq:affine-origin-a}
\end{equation}
Every solution on this layer is
\begin{equation}
 \widetilde d_r=\widetilde d_0+sr,
 \qquad
 u_r=\widetilde u_0+ar,
 \label{eq:affine-two-forms}
\end{equation}
and the determinant is the original fixed shift:
\begin{equation}
 \boxed{
 s\widetilde u_0-a\widetilde d_0=h_0\ne0.}
 \label{eq:affine-fixed-determinant}
\end{equation}

\begin{proposition}[Literal affine-column interface]
\label{prop:affine-column-interface}
For each fixed opposite row, polarization, and layer
\((\ell,j,s,v)\), the inner arithmetic sequence produced by the
actual mask has the form
\begin{equation}
 \boxed{
 \cC_{\ell,j,s,v}(I)
 =
 \sum_{r\in I}
 \mu(\widetilde d_0+sr)
 \mu(\widetilde u_0+ar)
 \one_{\left(v,\,
   (\widetilde d_0+sr)(\widetilde u_0+ar)\right)=1}
 \rho(r)W(r),}
 \label{eq:affine-literal-correlation}
\end{equation}
where:
\begin{enumerate}[label=\textup{(\alph*)},leftmargin=2.3em]
 \item \(a=\ell jv\),
 \((a,s)=1\), \((as,h_0)=1\), and
 \(s\widetilde u_0-a\widetilde d_0=h_0\);
 \item \(I\) is the union of at most two intervals and
 \begin{equation}
 |I|\ll 1+\frac{D}{sv};
 \label{eq:affine-fiber-length}
 \end{equation}
 \item \(\rho\) is a bounded fixed-period factor and \(W\) retains
 the logarithmic and smooth physical weights;
 \item the outside coefficient retains the opposite row, the
 row-gcd divisor coefficient, and the remaining physical factors;
 \item summation of the divisor coefficients for a fixed opposite
 squarefree row costs \(X^{o(1)}\), as does expansion of the displayed
 local coprimality condition.
\end{enumerate}
The symmetric construction covers the second polarization.
\end{proposition}

\begin{proof}
This is the actual-mask column disintegration of TPC-33
\citep{WangTPC33}, with the shift renamed \(h_0\).  The congruence
preceding \eqref{eq:affine-origin-a} gives
\eqref{eq:affine-two-forms} and
\eqref{eq:affine-fixed-determinant}.  Sampling a row interval of
length \(O(D/v)\) in a residue class modulo \(s\) gives
\eqref{eq:affine-fiber-length}.  Exact sign fusion on squarefree
support supplies the two M\"obius factors.  The row-gcd projector is
expanded only after the opposite row is fixed, yielding the
\(X^{o(1)}\) divisor mass.  Exchanging the row labels gives the other
polarization.
\end{proof}

\begin{remark}[Column complexity is not matrix complexity]
\label{rem:affine-mask-complexity}
The \(X^{o(1)}\) divisor cost in
\cref{prop:affine-column-interface} is a statement for one fixed
opposite row.  It is not an \(X^{o(1)}\) projective decomposition of
the complete joint row-pair mask.  Returning the column estimates
through the outer row sum remains part of the physical gate.
\end{remark}

\subsection{Integral normal form}

The geometry in \eqref{eq:affine-fixed-determinant} has a useful
coordinate-free description.  Assume
\[
 a,s\in\N,\qquad (a,s)=1,\qquad(as,h_0)=1.
\]
Choose \(x,y\in\Z\) such that \(sx-ay=1\).

\begin{proposition}[Fixed-determinant normal form]
\label{prop:affine-normal-form}
Every integral solution of
\begin{equation}
 sU-aD=h_0
 \label{eq:affine-normal-equation}
\end{equation}
has a unique representation
\begin{equation}
 U_t=h_0x+at,\qquad
 D_t=h_0y+st,\qquad t\in\Z.
 \label{eq:affine-normal-forms}
\end{equation}
Furthermore,
\begin{equation}
 \gcd(D_t,U_t)
 =\gcd(D_t,h_0)
 =\gcd(U_t,h_0)
 =\gcd(h_0,t).
 \label{eq:affine-gcd-identities}
\end{equation}
Whenever \(D_t,U_t>0\), on the primitive layer \((t,h_0)=1\),
\begin{equation}
 \mu(D_t)\mu(U_t)=\mu(D_tU_t),
 \label{eq:affine-mobius-fusion}
\end{equation}
and the reducible quadratic \(D_tU_t\) has discriminant \(h_0^2\).
\end{proposition}

\begin{proof}
The matrix
\[
 \begin{pmatrix}x&a\\y&s\end{pmatrix}
\]
has determinant one.  Applying it to \((h_0,t)^{\mathsf T}\) gives
\eqref{eq:affine-normal-forms}; applying its integral inverse proves
uniqueness.  A unimodular map preserves the ideal generated by the
two coordinates, which gives
\(\gcd(D_t,U_t)=\gcd(h_0,t)\).  The remaining gcd identities follow
from \((as,h_0)=1\).  On the primitive layer the two forms are
coprime, proving \eqref{eq:affine-mobius-fusion}.  Finally,
\[
 D_tU_t
 =as t^2+h_0(sx+ay)t+h_0^2xy
\]
has discriminant
\[
 h_0^2\{(sx+ay)^2-4asxy\}
 =h_0^2(sx-ay)^2=h_0^2.
\]
\end{proof}

\begin{remark}[Structural identity, not cancellation]
The broad normal-form class contains the classical adjacent
correlation: with \(a=s=h_0=1\), \(x=1\), and \(y=0\), one obtains
\(\mu(t)\mu(t+1)\).  Therefore a theorem uniform over the entire
normal-form class would contain an open Chowla case.  This observation
does not prove that the narrower physical subfamily in
\eqref{eq:affine-literal-correlation} is equivalent to that classical
case.
\end{remark}

\subsection{The complete reassembly gate}

Let \(\theta\) collect the fixed opposite row, polarization,
\((\ell,j,s,v)\), interval component, residue factor, and smooth
separation index.  An exact use of
\cref{prop:affine-column-interface} must have the form
\begin{equation}
 Z_X
 =\sum_{\theta\in\Theta_X}
 W_{\theta,X}\cC_{\theta,X}
 +E_{\rm content},
 \label{eq:affine-global-interface}
\end{equation}
where \(\cC_{\theta,X}\) has
\eqref{eq:affine-literal-correlation} and
\begin{equation}
 E_{\rm content}=-\cS_{>C}^{\rm sh}.
 \label{eq:affine-content-error}
\end{equation}
This is the reassembly of the raw hard coefficient
\eqref{eq:literal-AC}.  The calibrated drift legs are not part of
\(A_C\), \(Z_X\), or \(D_X\).
Equation \eqref{eq:affine-global-interface} is an interface, not an
estimate.

\begin{proposition}[Sufficient uniform affine budget]
\label{prop:affine-sufficient-budget}
Suppose \eqref{eq:affine-global-interface} has been verified with all
native keys and both polarizations, and suppose that, for some
\(\eta_Z\ge0\),
\begin{align}
 \sum_{\theta\in\Theta_X}
 |W_{\theta,X}\cC_{\theta,X}|
 &\le X^{-\eta_Z+o(1)}Q^2,
 \label{eq:affine-correlation-budget}\\
 |E_{\rm content}|
 &\le X^{-\eta_Z+o(1)}Q^2.
 \label{eq:affine-error-budget}
\end{align}
Then
\begin{equation}
 |Z_X|\le X^{-\eta_Z+o(1)}Q^2.
 \label{eq:affine-zero-bound}
\end{equation}
\end{proposition}

\begin{proof}
Apply the triangle inequality once, after the exact global
reassembly \eqref{eq:affine-global-interface}, and use
\cref{eq:affine-correlation-budget,eq:affine-error-budget}.
\end{proof}

\begin{remark}[Why the sufficient condition is demanding]
The absolute sum in \eqref{eq:affine-correlation-budget} is not
asserted here.  It may discard useful outer cancellation and can cost
a fixed power.  A stronger future argument may instead estimate the
complete outer operator.  Either route must still cover growing
\(a,s\), the \(O(1)\)-length fibers allowed by
\eqref{eq:affine-fiber-length}, the actual mask, and the same fixed
\(h_0\).
\end{remark}

\begin{remark}[The calibrated drift is downstream]
\label{rem:affine-drift-downstream}
TPC-32 writes the complete calibrated cutoff difference as the raw
matched shell plus two explicitly bounded drift legs.  Those legs
must be restored when the raw estimate is returned to the calibrated
physical packet, and their loss belongs in the complete downstream
ledger.  Inserting them into
\eqref{eq:affine-global-interface}, however, would change its left
side from the raw \(Z_X\) defined in \eqref{eq:literal-Z-D} to a
different calibrated zero mode.  The two normalizations are therefore
kept separate.
\end{remark}
